% ============================================================
% f2015tjl_q10_views.tex —— 2015 天津卷（理科）第 10 题 三视图
% ------------------------------------------------------------
% 坐标口径：图内画布 pt，原点 = 源页基准点，x 右正、y 上正。
%   来源：.work/phase3-b11/figure-inventory.md（真实页 pt，原点左上、y 向下）
%   换算：X_tikz = x_page − 493.36 ；  Y_tikz = 573.70 − y_page
% 比例基准：源题「1」= 28.35 pt（清单 §3），全图尺寸标注均为其整数倍。
% 线宽：图形本体 0.398 pt（轮廓 / 分隔线 / 尺寸线 / 界限短线统一本档）。
% 箭头：TikZ arrows.meta 标准箭矢（> = Triangle）；尺寸线每段一枚实心三角，
%       箭尖落在段尾坐标（= 清单所列「起点→终点」的终点，指向最近的尺寸界限线）。
% 照源复刻（清单 §7 存疑）：
%   S-5 俯视图**无水平尺寸线**（正视图有「1 | 2 | 1」三条）——不补。
%   S-6 侧视图半径线本体止于 (196.09,159.24)，箭尖 (198.63,161.78) 至圆心 28.35 = 恰好 1 图；
%       「线段长」与「半径」口径不一致——照源，半径以尖端计。
% 要素：46 描边 path（正视 27 + 侧视 2 + 俯视 17）+ 23 填充箭头（14 + 1 + 8）
%       + 9 文字标签（清单 §3.5 B 表、§3.4）。
% ============================================================
\usetikzlibrary{arrows.meta}
\begin{tikzpicture}[
  x=1pt,y=1pt,                             % 画布单位 = 1 pt（本文件坐标一律为 pt）
  every node/.style={font=\fontsize{10.5pt}{12.6pt}\selectfont,inner sep=0pt},
  >={Triangle[length=3.60pt,width=3.36pt]},
  g/.style={line width=0.398pt},
]

  % =========================================================
  % 正视图（上左）
  % =========================================================
  % 轮廓（6 点六边形：中央 2 图×2 图矩形 + 左右各伸 1 图三角尖）
  \draw[g] (8.50,141.74) -- (36.85,113.39) -- (93.54,113.39) -- (121.89,141.74)
           -- (93.54,170.08) -- (36.85,170.08) -- cycle;
  % 内部分隔竖线
  \draw[g] (36.85,113.39) -- (36.85,170.08);
  \draw[g] (93.54,113.39) -- (93.54,170.08);
  % 尺寸箭头：竖直 · 右外 @ X=127.56（每组「1」上下各一枚，箭尖朝外）
  \draw[g,->] (127.56,121.89) -- (127.56,116.98);
  \draw[g,->] (127.56,133.23) -- (127.56,138.15);
  \draw[g,->] (127.56,150.24) -- (127.56,145.32);
  \draw[g,->] (127.56,161.58) -- (127.56,166.50);
  % 尺寸箭头：竖直 · 左外 @ X=2.83
  \draw[g,->] (2.83,121.89) -- (2.83,116.98);
  \draw[g,->] (2.83,133.23) -- (2.83,138.15);
  \draw[g,->] (2.83,150.24) -- (2.83,145.32);
  \draw[g,->] (2.83,161.58) -- (2.83,166.50);
  % 尺寸箭头：水平 · 下方 @ Y=107.72（三段读数 1 | 2 | 1）
  \draw[g,->] (17.00,107.72) -- (12.09,107.72);
  \draw[g,->] (28.34,107.72) -- (33.26,107.72);
  \draw[g,->] (59.52,107.72) -- (40.43,107.72);
  \draw[g,->] (70.86,107.72) -- (89.95,107.72);
  \draw[g,->] (102.04,107.72) -- (97.13,107.72);
  \draw[g,->] (113.38,107.72) -- (118.30,107.72);
  % 尺寸界限短线：左外 3 条 / 右外 3 条
  \draw[g] (5.67,113.39) -- (0.00,113.39);
  \draw[g] (5.67,141.74) -- (0.00,141.74);
  \draw[g] (5.67,170.08) -- (0.00,170.08);
  \draw[g] (124.72,113.39) -- (130.39,113.39);
  \draw[g] (124.72,141.74) -- (130.39,141.74);
  \draw[g] (124.72,170.08) -- (130.39,170.08);
  % 尺寸界限短线：下方竖短线 4 条 @ 轮廓四角 x
  \draw[g] (8.50,110.55)  -- (8.50,104.89);
  \draw[g] (36.85,110.55) -- (36.85,104.89);
  \draw[g] (93.54,110.55) -- (93.54,104.89);
  \draw[g] (121.89,110.55) -- (121.89,104.89);
  % 读数与图名
  \node at (22.64,107.00)  {1};
  \node at (65.24,107.00)  {2};
  \node at (107.74,107.00) {1};
  \node at (2.84,155.20)   {1};
  \node at (2.84,126.80)   {1};
  \node at (127.54,155.20) {1};
  \node at (127.54,126.80) {1};
  \node at (65.24,91.30)   {正视图};

  % =========================================================
  % 侧视图（上右）：半径 1 图圆 + 半径引导线
  % =========================================================
  \draw[g] (178.58,141.74) circle [radius=28.35pt];
  \draw[g,->] (178.58,141.74) -- (198.63,161.78);
  \node at (192.74,143.80) {1};
  \node at (178.54,91.30)  {侧视图};

  % =========================================================
  % 俯视图（下左）：与正视图同形同尺寸，沿 y 下移 113.39
  % =========================================================
  \draw[g] (8.50,28.35) -- (36.85,0.00) -- (93.54,0.00) -- (121.89,28.35)
           -- (93.54,56.70) -- (36.85,56.70) -- cycle;
  \draw[g] (36.85,0.00) -- (36.85,56.70);
  \draw[g] (93.54,0.00) -- (93.54,56.70);
  % 尺寸箭头：竖直 · 右外 @ X=127.56
  \draw[g,->] (127.56,8.51)  -- (127.56,3.59);
  \draw[g,->] (127.56,19.84) -- (127.56,24.76);
  \draw[g,->] (127.56,36.85) -- (127.56,31.93);
  \draw[g,->] (127.56,48.19) -- (127.56,53.11);
  % 尺寸箭头：竖直 · 左外 @ X=2.83
  \draw[g,->] (2.83,8.51)  -- (2.83,3.59);
  \draw[g,->] (2.83,19.84) -- (2.83,24.76);
  \draw[g,->] (2.83,36.85) -- (2.83,31.93);
  \draw[g,->] (2.83,48.19) -- (2.83,53.11);
  % 尺寸界限短线：左外 3 条 / 右外 3 条（照源无水平尺寸线，S-5）
  \draw[g] (5.67,0.00)  -- (0.00,0.00);
  \draw[g] (5.67,28.35) -- (0.00,28.35);
  \draw[g] (5.67,56.70) -- (0.00,56.70);
  \draw[g] (124.72,0.00)  -- (130.39,0.00);
  \draw[g] (124.72,28.35) -- (130.39,28.35);
  \draw[g] (124.72,56.70) -- (130.39,56.70);
  % 读数与图名
  \node at (2.84,41.80)   {1};
  \node at (2.84,13.40)   {1};
  \node at (127.54,41.80) {1};
  \node at (127.54,13.40) {1};
  \node at (65.24,-13.50) {俯视图};

\end{tikzpicture}
